%=============================================================================
%  10_prediction.tex  -  Chapter 9: The Prediction Engine
%=============================================================================
\chapter{The Prediction Engine}
\label{ch:predict}

The prediction skill lets an analyst go beyond description into lightweight
modelling: ``can this dataset predict X?'' It is a hybrid pipeline --- the
language model reasons about \emph{what} to model and \emph{which} algorithm to
use, while scikit-learn does the actual training and evaluation. The model
advises; deterministic code decides and measures.

\section{Pipeline stages}
\begin{figure}[h]
\centering
\begin{tikzpicture}[node distance=4mm and 7mm,
  s/.style={draw=uaIndigo,fill=uaIndigo!8,rounded corners=3pt,align=center,
    font=\footnotesize,inner sep=6pt,minimum height=12mm,text width=2.5cm},
  a/.style={-{Stealth[length=2.4mm]},uaSlate,thick}]
  \node[s] (t) {1. Target\\selection\\\scriptsize (LLM + schema)};
  \node[s,right=of t] (k) {2. Task type\\\scriptsize classif. / regr.};
  \node[s,right=of k] (m) {3. Model\\choice\\\scriptsize (LLM + schema)};
  \node[s,right=of m] (p) {4. Preprocess\\\scriptsize impute / encode / scale};
  \node[s,right=of p] (e) {5. Train +\\evaluate};
  \draw[a](t)--(k); \draw[a](k)--(m); \draw[a](m)--(p); \draw[a](p)--(e);
\end{tikzpicture}
\caption{The five-stage prediction pipeline, blending LLM reasoning (stages 1--3)
with scikit-learn execution (stages 4--5).}
\label{fig:predict}
\end{figure}

\section{Target and task selection}
The first decision --- which column to predict --- is made by the model, but
inside a schema. A generated JSON schema constrains the model to choose a
\code{target\_column} from the real column list, give a \code{reason}, and
declare the \code{task\_type} as either \code{classification} or
\code{regression}. Structured output turns this from a free-text guess into a
validated selection.

\begin{lstlisting}[style=py,caption={The target-selection schema (abridged).}]
{
  "target_column": {"type": "string", "enum": columns},
  "reason": {"type": "string"},
  "task_type": {"type": "string", "enum": ["classification", "regression"]}
}  # required: target_column, reason, task_type
\end{lstlisting}

\section{Model selection}
Given the task type, a second schema constrains the model to pick an appropriate
scikit-learn algorithm and its hyperparameters from a curated list. Offering the
model a fixed menu --- rather than free choice --- guarantees that whatever it
returns can actually be instantiated and trained.

\begin{table}[h]
\centering
\caption{Candidate models offered to the LLM per task type.}
\small
\begin{tabularx}{\linewidth}{@{}L{3.0cm} X@{}}
\toprule
\textbf{Task} & \textbf{Candidate estimators} \\
\midrule
Classification & \code{logistic\_regression}, \code{random\_forest\_classifier},
\code{gradient\_boosting\_classifier}, \code{decision\_tree\_classifier},
\code{svm\_classifier}, \code{knn\_classifier} \\
Regression & \code{linear\_regression}, \code{ridge}, \code{lasso},
\code{random\_forest\_regressor}, \code{gradient\_boosting\_regressor},
\code{decision\_tree\_regressor}, \code{svr}, \code{knn\_regressor} \\
\bottomrule
\end{tabularx}
\end{table}

\section{Preprocessing}
Before training, the pipeline prepares features deterministically with
scikit-learn utilities: missing values are imputed (\code{SimpleImputer}),
categorical features are label-encoded (\code{LabelEncoder}), and numeric
features are standardized (\code{StandardScaler}). The data is split into train
and test partitions with \code{train\_test\_split}. None of this touches the
language model --- it is ordinary, reproducible feature engineering.

\begin{lstlisting}[style=py,caption={The scikit-learn building blocks used by the pipeline.}]
from sklearn.model_selection import train_test_split
from sklearn.preprocessing import LabelEncoder, StandardScaler
from sklearn.impute import SimpleImputer
from sklearn.linear_model import LogisticRegression, LinearRegression, Ridge, Lasso
from sklearn.ensemble import (RandomForestClassifier, RandomForestRegressor,
                              GradientBoostingClassifier, GradientBoostingRegressor)
from sklearn.tree import DecisionTreeClassifier, DecisionTreeRegressor
from sklearn.svm import SVC, SVR
from sklearn.neighbors import KNeighborsClassifier, KNeighborsRegressor
\end{lstlisting}

\section{Training and evaluation}
The chosen estimator is fit on the training partition and evaluated on the held-
out test set with task-appropriate metrics: accuracy, precision, recall, F1, and
a classification report for classification; $R^2$ and error metrics for
regression. The results --- the selected target and its rationale, the chosen
model, and the metrics --- are streamed back so the user sees not just a number
but \emph{why} this model was chosen and how well it did.

\begin{table}[h]
\centering
\caption{Evaluation metrics by task type.}
\small
\begin{tabularx}{\linewidth}{@{}L{3.0cm} X@{}}
\toprule
\textbf{Task} & \textbf{Metrics reported} \\
\midrule
Classification & Accuracy, precision, recall, F1, and a full classification
report. \\
Regression & $R^2$ score and error measures. \\
\bottomrule
\end{tabularx}
\end{table}

\begin{uanote}[Bounded autonomy]
The model is given real agency --- it selects the target and the algorithm ---
but only within schemas and a fixed estimator menu. This is the project's
signature pattern applied to machine learning: \emph{the model proposes within a
contract; deterministic code disposes and measures.}
\end{uanote}
